\begin{singlespace}
\begin{abstract}
A payment can go through even when the person making it needs someone else's account. We use a referral survey of 423 migrant women in Colombia, written syntheses of interviews and focus groups, and institutional evidence to study what digital payment counts miss. The survey screened for international remittance receipt by respondents or their households in the previous year. Among women answering each question, 89.1 percent reported sending or receiving money digitally within Colombia in the past 60 days, and 40.0 percent reported keeping or using some incoming remittance proceeds digitally. In separate logit models, reporting both a phone line and proof of address in one's own name rather than one is positively associated with retaining or using remittances digitally. A higher score combining reported transfer ability and confidence in avoiding fraud is positively associated with domestic use. Women described paying through relatives' accounts and relying on intermediaries for delivery to family abroad. Payment counts alone cannot tell us who controls the account, when the recipient can use the funds, or where the user can turn if a payment fails.
\end{abstract}
\small
\noindent\textbf{Keywords:} digital payments; remittances; migrant women; financial inclusion; digital public infrastructure; Colombia.
\par\noindent\textbf{JEL codes:} O33, G21, F24, J15, D14.
\end{singlespace}
